\section{CONCLUSIONS AND RECOMMENDATIONS}

\subsection{Conclusions}
\hspace{1.27cm}
This thesis demonstrated the complete design, development, and experimental validation of a vision-based Delta robot pick-and-place system using ROS 2 intended for educational and light industrial applications in Cambodia. The following conclusions are drawn from the work carried out:

\begin{itemize}
    \item A fully functional Delta robot prototype was successfully designed and fabricated using RobStride RS-00 quasi-direct-drive actuators, aluminum structural profiles, and 3D-printed PLA linkages, with geometric parameters $f=157$\,mm, $e=35$\,mm, $r_f=200$\,mm, $r_e=400$\,mm, demonstrating that capable robotic hardware can be assembled from accessible components at reduced cost.

    \item The forward and inverse kinematic models of the Delta robot were derived analytically following the Trossen Robotics geometric formulation and validated through physical experiments, achieving FK position errors consistently below 0.5\,mm across the full workspace. All joint angle solutions were constrained to $\theta_i \geq 0$ to ensure physical reachability.

    \item A complete ROS\,2 control architecture was implemented, comprising five custom packages (\texttt{delta\_common}, \texttt{delta\_motor\_controller}, \texttt{delta\_camera\_system}, \texttt{delta\_main\_app}, and \texttt{delta\_can\_driver}), communicating via CAN bus at 1\,Mbit/s and coordinated by an 8-state finite state machine for pick-and-place task sequencing.

    \item A full vision pipeline was implemented and validated, encompassing Intel RealSense D455 RGB-D capture, HSV colour-space object detection (confidence 1.00), pinhole back-projection for 3D localisation using \texttt{FAKE\_DEPTH\_M = 0.650}\,m, and a calibrated rigid-body camera-to-robot transform ($T_x=-10.0$, $T_y=+235.0$, $T_z=-20.0$\,mm) verified through a projected workspace overlay with $(0,0)$ crosshair marker.

    \item A conveyor belt prediction algorithm was developed using a triangular motion profile ($t = 2\sqrt{d/a_{\max}}$), with belt speed physically measured at 12.85\,mm/s, yielding a worst-case belt offset of 8.3\,mm at medium speed preset.

    \item A visual end-effector feedforward correction loop was implemented, detecting a laser dot EE marker via HSV masking in the raw camera frame and iteratively correcting the approach position before descent. The X-axis correction converged to below the 3\,mm threshold (\texttt{EE\_CORRECTION\_THRESH\_MM = 3.0}) within 2--4 iterations for stationary objects.

    \item The robot successfully executed the complete 8-state pick-and-place sequence (IDLE $\to$ APPROACHING $\to$ PICKING $\to$ GRASPING $\to$ LIFTING $\to$ TRANSPORTING $\to$ DROPPING $\to$ RESETTING $\to$ IDLE) with FK errors below 0.5\,mm on all moves, demonstrating reliable end-to-end operation.
\end{itemize}

\subsection{Recommendations}
\hspace{1.27cm}
To further enhance the system and broaden its applications, the following recommendations are proposed for future work:

\begin{itemize}
    \item \textbf{Non-linear error map calibration:} Complete the grid calibration sweep using \\ \texttt{measure\_error\_grid.py}, which commands the robot to 81 grid points and uses camera-tracked EE position to build a bilinearly interpolated error correction map. Enabling \texttt{ERROR\_MAP\_ENABLE = True} is expected to reduce residual pick errors to below 2\,mm across the full workspace.

    \item \textbf{Y-axis visual correction for belt picks:} The current correction loop fixes X-axis error only. For stationary objects, full XY correction already converges reliably. For belt picks, a predictive Y-correction strategy that accounts for object motion during iteration should be developed to extend correction to both axes.

    \item \textbf{Gripper force and contact feedback:} Adding a force sensor or contact switch to the pneumatic gripper would enable pick success verification, allowing the FSM to detect failed grasps and retry, significantly improving system reliability.

    \item \textbf{Structural simulation:} Future iterations should include finite element analysis (FEA) using software such as ANSYS or SolidWorks Simulation to evaluate stress, strain, and safety factors of the mechanical components under dynamic loading at higher speed presets.

    \item \textbf{Multi-class detection and sorting:} Training the YOLO detector on multiple object classes and implementing a routing logic to multiple drop positions would enable automated sorting applications relevant to food processing and electronics assembly.

    \item \textbf{Performance benchmarking:} A systematic evaluation of cycle time, pick success rate, and repeatability across multiple trials and positions should be conducted to provide quantitative performance characterisation for publication and comparison with related work.

    \item \textbf{Higher speed preset validation:} The system was validated at the medium preset ($v_{\max}=2000$\,mm/s). Testing at the maximum preset ($v_{\max}=4000$\,mm/s) with safety guards and FK error monitoring would characterise the upper performance limit of the hardware.

    \item \textbf{Enhanced user interface:} A real-time RViz2 or custom GUI with workspace visualisation, live FK position display, and calibration tools would improve usability for both educational demonstrations and industrial deployment.
\end{itemize}
